\scanpage{84}
% archive foliation: 86
\fbox{\unsure{Modulis|Moduli} $M_{0,4}$ et l'octaèdre}

$I$ ens de card. 4 d'où
\[
J=\unsure{P|\mathcal P}_{2,2}(I),\quad V(J)\ \text{plan vectoriel sur }\mathbb F_2\text{ avec }V(J)^*\simeq J,
\]
$V(J)\subset\unsure{\mathfrak S}_I$ formé de l'identité et des ``bitranspositions''.

On s'intéresse aux droites projectives (sur $\mathbb C$ disons --- \unsure{ça}
marcherait sur tt corps $k$ \inserted{de car. $\ne2$}) munies de $I\hookrightarrow X$.
(plus précisément $I_{\mathbb C}\hookrightarrow X$, ou $I\longrightarrow X(\mathbb C)$ \ldots).
L'opération de $V(\unsure{I|J})$ sur $I$ se prolonge en une opération de
$V(J)$ sur $X$
\[
V(J)\hookrightarrow G_X=\operatorname{Aut}(X)
\]
Soit
\[
\begin{gathered}
S=\text{ens des pts dans }X\text{ qui sont fixes par un}\\
\text{élément de }J=V(J)^*\\
\operatorname{card}S=6
\end{gathered}
\]
\margin{\uline{\unsure{Commençons}} \uline{par \unsure{ça}}}
\figure{\begin{tikzpicture}\draw (0,0)--(0,1.1);\end{tikzpicture}}
En fait, on trouve que \deleted{\unsure{pour les}}\inserted{les} opérations \inserted{fidèles} d'un gr. abstrait $V$
sur un $X$, i.e. les plongements $V(\unsure{J})\hookrightarrow G_X$, sont tous
équivalents (modulo autom. intérieurs de $G_X$)
et sont (pour une base fixée \inserted{$e_0,e_1$}) équivalentes à l'opération sur $\mathbb P^1$
\diagram{\begin{array}{rcll}
u(e_0)&=(z\longmapsto-z)=\begin{pmatrix}-1&0\\0&1\end{pmatrix}& &\text{pts fixes }0,\infty\\[3pt]
u(e_1)&=(z\longmapsto1/z)=\begin{pmatrix}0&1\\1&0\end{pmatrix}& &\text{pts fixes }\pm1\\[3pt]
u(e_0+e_1)&=(z\longmapsto-1/z)=\begin{pmatrix}0&-1\\1&0\end{pmatrix}& &\text{pts fixes }\pm i
\end{array}}
Les \deleted{\unsure{quatre groupes}}\inserted{sous-groupes} isom. à $D_2$ de $G_X$ correspondent 1-1 aux
parties $S$ de cardinal 6 de $X$, isom. à $(\mathbb P_1,\{0,\infty,\pm1,\pm i\})$.
\margin{parties octaédriques d'une droite projective}
Le normalisateur $\Gamma$ d'un tel \deleted{$\unsure{D_2}$} sous-groupe $V$ \inserted{dans $G_X$} est isomorphe
au groupe \deleted{des rotations} $\unsure{\mathfrak S}_4$ (de façon non canonique), plus
précisément, si $V=V(J)$ avec $J=\unsure{P|\mathcal P}(I)$ \inserted{i.e. si $I$ est un $V$-torseur}, on trouve \deleted{des} iso
de suites exactes
\diagram{\begin{tikzcd}[ampersand replacement=\&,column sep=small]
\&\& \deleted{\unsure{\operatorname{Aut}(I)}} \arrow[d,equal] \&\& \\
\unsure{1} \arrow[r] \& V \arrow[r] \arrow[d,equal] \& \unsure{\mathfrak S}_I \arrow[r] \arrow[d] \& \operatorname{Aut}(V) \arrow[r,dash] \arrow[d,equal] \& 1 \\
1 \arrow[r] \& V \arrow[r] \& \Gamma \arrow[r] \& \operatorname{Aut}(V) \arrow[r,dash] \& 1
\end{tikzcd}}
défini modulo autom. intérieurs par un élément de $V$.
Quand $I$ est le torseur trivial, \unsure{se donner|on donne} un tel \deleted{iso}\inserted{iso d'ext.}
